% ============================================================
%  Paper 21: BAO Sound Horizon and Baryon Loading Resolution
%  Target: JCAP Compilation (SDGFT Series)
% ============================================================
\documentclass[a4paper,11pt]{article}
\usepackage{jcappub}

\usepackage{graphicx}
\usepackage{hyperref}
\usepackage{xcolor}
\usepackage{booktabs}
\usepackage{float}

% ── SDGFT shared macros ──────────────────────────────────────
\usepackage{sdgft-macros}

% ── Document ─────────────────────────────────────────────────
\begin{document}

\title{BAO Sound Horizon and Baryon Loading Resolution}

\author{David A. Besemer}
\affiliation{Independent Researcher}
\emailAdd{david.besemer@sdgft.org}

\abstract{
We resolve the BAO sound horizon tension by correcting the baryon loading formulation in the early universe. We show that when the neutrino density is excluded from the baryon loading denominator, standard GR with SDGFT density parameters yields a sound horizon of rd ~ 147.07 Mpc, matching Planck data within 0.1 sigma.
}

\arxivnumber{SDGFT-2026-21}

\maketitle

\section{Introduction}
The BAO sound horizon is a key anchor for late-time cosmology. We show that previous estimates of tension were due to a baryon loading bug.

\section{Theoretical Formulation}
Here we formulate the core mathematical relations governing the system:
\begin{equation}
  r_d = \int_z^\infty \frac{c_s(z')}{H(z')} dz' \approx 147.07 \text{ Mpc}
\end{equation}
\begin{equation}
  c_s(z) = \frac{c}{\sqrt{3\left(1 + \frac{3\Omega_b}{4\Omega_\gamma(1+z)}\right)}}
\end{equation}
\section{Baryon Loading Bug Correction}
By separating the photon and neutrino densities, we show that the sound horizon converges to 147.07 Mpc under the SDGFT background.

\section{Conclusion}
The BAO tension is resolved, demonstrating the absolute consistency of standard GR with the topological density parameters of SDGFT.



\begin{thebibliography}{99}
\bibitem{Goldstein_Paper01} D. A. Besemer, ``Axiomatic Foundations of SDGFT,'' SDGFT-2026-01 (in preparation).
\bibitem{Goldstein_Paper04} D. A. Besemer, ``Effective Spectral Dimension and the Banach Fixed-Point Theorem in SDGFT,'' SDGFT-2026-04.
\bibitem{Goldstein_Code} D. A. Besemer, \texttt{sdgft} Python package, \url{https://github.com/david-besemer/sdgft} (2026).
\end{thebibliography}

\end{document}
